
%%%%%%%%%%%%%%%%%%%%%%%%%%%%%%%%%%%%%%%%%%%%%%%%%%%%%%%%%%%%%
\subsubsection{模型建立}

针对问题四，本节建立以边际价值为核心的数据建议排序框架。每条建议按两个维度评估：对模型关键薄弱环节的直接修正作用，以及采集可行性与成本。据此将建议分为高、中、低三个优先级，共九条。

\textbf{步骤1：高优先级建议（3条）}

（1）每日实际库存与缺货记录。当前数据仅含销售流水，无法区分"零销量=无需求"和"零销量=缺货"。87个单品销售天数不足30天、5个单品从未销售——这些究竟是需求不足还是供应缺失，现有数据无法判断。采集每日库存数据可将需求估计从"销售量"修正为"销售量+缺货量"，直接改善Prophet预测和报童模型的需求输入。

（2）折扣原因分类码。折扣回收率$k$是模型最敏感参数——问题二利润在$k\in[0.50,0.75]$区间内波动44\%。当前$k=0.66$为所有打折交易的全局中位数，但打折原因混杂：运损清仓（品相变差被迫打折）vs库存过剩清仓（订货过多主动打折）。记录折扣原因码可分别估计运损折扣率和清仓折扣率，显著提高报童模型参数精度。

（3）天气数据（温度、湿度、降水量）。Prophet残差仍有显著的自相关结构（Ljung-Box $p\approx 0$），表明存在模型未捕捉的外部驱动因素。天气是蔬菜需求最可能的遗漏变量——高温加速蔬菜变质、暴雨减少客流、季节性温湿度影响特定品类需求。将天气数据作为Prophet的外生回归项引入，可降低残差自相关，提高预测精度。

\textbf{步骤2：中优先级建议（3条）}

（4）单品上架与下架时间表，可区分季节性商品（特定季节才上架）和低销量常驻商品，使Q3候选单品池不被动依赖最近一周的销售快照。（5）促销活动日历，Prophet的内置中国节假日仅覆盖法定假日，不含商超自发的促销活动（如会员日折扣），作为额外虚拟变量可分离促销效应和自然需求。（6）竞争对手价格数据，当前从自身数据的log-log回归估计价格弹性面临遗漏变量偏误——品类销售量的变化可能部分来自竞争对手的价格调整，竞争对手数据可改善交叉弹性估计。

\textbf{步骤3：低优先级建议（3条）}

（7）消费者偏好与满意度调查，帮助识别上升或衰落中的单品，改善Q3单品选择的前瞻性。（8）货架空间与陈列数据，可将Q3的定性空间约束转化为定量模型输入，替代当前以SKU数量间接表达的空间限制。（9）产地与供应链数据，附件1中部分单品名称已含供应来源编号（如数字编号表示不同供应来源），完整提供有助于批发价波动预测和供应风险评估。

\subsubsection{模型求解}

基于上述分析框架，将九条建议按优先级汇总如表\ref{tab:q4summary}所示。

\begin{table}[H]
  \centering
  \caption{数据采集建议汇总}
  \label{tab:q4summary}
  \begin{tabularx}{\textwidth}{CCC}
    \toprule
    优先级 & 建议数据 & 对应建模薄弱环节 \\
    \midrule
    高 & 每日库存与缺货记录 & 需求-销量识别偏差（Q2, Q3） \\
    高 & 折扣原因分类码 & $k$参数精度（Q2, Q3） \\
    高 & 天气数据 & Prophet残差结构（Q2） \\
    中 & 单品上下架时间表 & 候选单品池时效性（Q1, Q3） \\
    中 & 促销活动日历 & 需求异常值识别（Q2） \\
    中 & 竞争对手价格数据 & 弹性估计遗漏变量（Q2） \\
    低 & 消费者偏好调查 & 单品选择前瞻性（Q3） \\
    低 & 货架空间与陈列数据 & 空间约束定量化（Q3） \\
    低 & 产地与供应链数据 & 批发价波动预测（Q2） \\
    \bottomrule
  \end{tabularx}
\end{table}

由表\ref{tab:q4summary}可知，高优先级的三类数据（库存记录、折扣原因码、天气数据）对应模型中三个最关键的薄弱环节，且采集成本低（系统自动记录或从公开API获取）。中优先级的数据对模型精度有可量化改善但需要额外投入。低优先级数据有益但对模型结构的改变大于对参数的修正。

综上所述，本文基于Q1--Q3建模过程中实际暴露的数据局限性，按边际价值排序提出了九条数据采集建议。建议的优先级反映了每条数据对模型关键薄弱环节的直接修正作用和采集可行性，为商超的数据采集决策提供了量化依据。
